\documentclass[12pt]{article}
\usepackage[left=1cm,right=1cm,top=0.5cm,bottom=1.5cm,bindingoffset=0cm]{geometry}
\usepackage[T1,T2A]{fontenc}
\usepackage[utf8]{inputenc}
\usepackage[english,russian]{babel}
% Пакет, необходимый для корректной работы babel
\usepackage{csquotes}
\usepackage{graphicx}
% Описание пути к изображениям.
\graphicspath{ {./pics} }
\usepackage{amsmath}
\usepackage{amsfonts}
\usepackage{amssymb}
\usepackage{amsthm}
\usepackage{mathtools}
% Пакет для удобной работы с системами уравнений
\usepackage{systeme}
% Пакет, позволяющий рисовать красивые картинки прямо внутри пдф файла
\usepackage{tikz}
\usepackage{tikz-3dplot}
\usepackage{hyperref}
\usepackage{xcolor}
\definecolor{linkcolor}{HTML}{225ae2}
\definecolor{urlcolor}{HTML}{225ae2}
\hypersetup{
    pdfstartview=FitH, 
    linkcolor=linkcolor,
    urlcolor=urlcolor,
    colorlinks=true
}
% Выделение ссылок
\newcommand\tab[1][.5cm]{\hspace*{#1}}
%отступ после начала окружения
\usepackage{bbm}
\usepackage[normalem]{ulem}
\theoremstyle{plain}
\newtheorem{theorem}{Теорема}[section]
\newtheorem{consequence}{Следствие}[theorem]
\newtheorem{subtheorem}{Утверждение}[section]
\newtheorem{algorithm}{Алгоритм}[section]
\newtheorem{axiom}{Аксиома}[section]
\newtheorem*{exercise}{Упражнение}
\newtheorem{propertes}{Свойства}
\newtheorem{lemma}{Лемма}[section]
\newtheorem{example}{Пример}
\newtheorem{examples}{Примеры}
\newtheorem{statement}{Формулировка}
% Окружение типа ``определение''
\newtheorem{definition}{Определение}[section]
\newtheorem{remark}{Замечание}
% Замечение без нумерации
\renewcommand\qedsymbol{$\blacksquare$}
\newcommand{\mP}{\mathbf{P}}
\newcommand{\mE}{\mathbf{E}}
\newcommand{\mD}{\mathbf{D}}
\newcommand{\skob}[1]{\left(#1\right)}
\newcommand{\fig}[1]{\left\{#1\right\}}
\newcommand{\br}[1]{\left[#1\right]}
\newcommand{\R}{\mathbb R}
\newcommand{\Q}{\mathbb Q}
\newcommand{\Z}{\mathbb Z}
\newcommand{\N}{\mathbb N}
\newcommand{\CC}{\mathbb C}
\newcommand{\acts}{\curvearrowright}
\newcommand{\F}{\mathbb F}
\newcommand{\aug}{\fboxsep=-\fboxrule\!\!\!\fbox{\strut}\!\!\!}
\newcommand{\sgn}{\operatorname{sgn}}
\newcommand{\id}{\mathrm{id}}
\renewcommand{\phi}{\varphi}
\renewcommand{\epsilon}{\varepsilon}
\newcommand{\E}{\mathcal E}
\begin{document}
    \tableofcontents
    \fontsize{14pt}{20pt}\selectfont
    \newpage
    \section{Тензоры и тензорные поля}
    \begin{definition}
        Пусть $P\in M^n$ --- точка на многообразии, тогда тензором типа $(p, \ q)$ ранга $p+q$ называется соответствие $(x^1, \ \ldots \ x^n) \ \longrightarrow \ T^{i_1,\ldots,i_p}_{j_1,\ldots,j_q}$ такое, что при замене базиса
        $$T^{i_1',\ldots,i_p'}_{j_1',\ldots,j_q'}=T^{i_1,\ldots,i_p}_{j_1,\ldots,j_q}\cdot\frac{\partial x^{i_1'}}{\partial x^{i_1}}\cdot\ldots\cdot\frac{\partial x^{i_p'}}{\partial x^{i_p}}\cdot\frac{\partial x^{j_1}}{\partial x^{j_1'}}\cdot\ldots\cdot\frac{\partial x^{j_q}}{\partial x^{j_q'}}$$
        Мультииндексная запись
        $$T_{J_q'}^{I_p'}=T_{J_q}^{I_p}\cdot\frac{\partial x^{I_p'}}{\partial x^{I_p}}\cdot\frac{\partial x^{J_q}}{\partial x^{J_q'}}$$
    \end{definition}
    \begin{examples}\tab
        \begin{enumerate}
            \item $T_2^1$, $T^{i'}_{j'k'}=T_{jk}^i\cdot\frac{\partial x^{i'}}{\partial x^{i}}\cdot\frac{\partial x^{j}}{\partial x^{j'}}\cdot\frac{\partial x^{k}}{\partial x^{k'}}$
            \item $T_0^1$, $(x^1, \ \ldots, \ x^n) \ \longrightarrow \ x^i$ вектор
            \item $T_1^0$, $(x^1, \ \ldots, \ x^n) \ \longrightarrow \ x_i$ ковектор (линейный функционал)
            \item $T_2^0$, $(x^1, \ \ldots, \ x^n) \ \longrightarrow \ (a^i_j)$ билинейная форма
            $$T_{i'j'}=T_{ij}\cdot\frac{\partial x^{i}}{\partial x^{i'}}\cdot\frac{\partial x^{j}}{\partial x^{j'}}$$
        \end{enumerate}
    \end{examples}
    \begin{lemma}
        Пусть $(x^1, \ \ldots, \ x^n)$ --- система координат в окрестности точки $P\in M$. Рассмотрим соответствие $(x^1, \ \ldots, \ x^n) \ \longrightarrow \ T^{I_p}_{J_q}$. По тензорному закону\\
        $(x^{1'}, \ \ldots, \ x^{n'})$ $\longrightarrow$ $T^{I_p'}_{J_q'}=T^{I_p}_{J_q}\cdot\frac{\partial x^{I_p'}}{\partial x^{I_p}}\cdot\frac{\partial x^{J_q}}{\partial x^{J_q'}}$. Тогда это соответствие тензор $T_q^p$.
    \end{lemma}
    \begin{proof}
        Рассмотрим систему координат $(x^{1''}, \ \ldots, \ x^{n''})$, по тензорному закону
        $$(x^{1''}, \ \ldots, \ x^{n''}) \ \longrightarrow \ T^{I_p''}_{J_q''}=T^{I_p}_{J_q}\cdot\frac{\partial x^{I_p''}}{\partial x^{I_p}}\cdot\frac{\partial x^{J_q}}{\partial x^{J_q''}}$$
        Заметим, что $T^{I_p'}_{J_q'}=T^{I_p}_{J_q}\cdot\frac{\partial x^{I_p'}}{\partial x^{I_p}}\cdot\frac{\partial x^{J_q}}{\partial x^{J_q'}} \ \Longleftrightarrow \ T^{I_p}_{J_q}=T^{I_p'}_{J_q'}\cdot\frac{\partial x^{I_p}}{\partial x^{I_p'}}\cdot\frac{\partial x^{J_q'}}{\partial x^{J_q}}$. Тогда
        $$T^{I_p''}_{J_q''}=T^{I_p'}_{J_q'}\cdot\frac{\partial x^{I_p}}{\partial x^{I_p'}}\cdot\frac{\partial x^{J_q'}}{\partial x^{J_q}}\cdot\frac{\partial x^{I_p''}}{\partial x^{I_p}}\cdot\frac{\partial x^{J_q}}{\partial x^{J_q''}}=T^{I_p'}_{J_q'}\cdot\frac{\partial x^{I_p''}}{\partial x^{I_p'}}\cdot\frac{\partial x^{J_q'}}{\partial x^{J_q''}}$$
    \end{proof}
    \begin{definition}
        Тензорным полем называется соответствие, сопоставляющее каждой точке поверхности тензор.
    \end{definition}
    \begin{definition}
        Пусть $M$ --- поверхность, $T_PM=V$ --- линейное пространство в точке $P\in M$, $V^*=T_P^*M$ --- пространство линейных функций. Пусть $(x^1, \ \ldots, \ x^n)$ --- базис $V$, $(dx^1, \ \ldots, \ dx^n)$ --- базис $V^*$. Рассмотрим $\underbrace{V\times \dots\times V}_p \times \underbrace{V^*\times \dots\times V^*}_q$ и полилинейное отображение $f$, тогда тензором $T^p_q$ называется $f:\underbrace{V\times \dots\times V}_p \times \underbrace{V^*\times \dots\times V^*}_q$.
    \end{definition}
    \begin{subtheorem}
        Определения тензора равносильны.
    \end{subtheorem}
    \begin{proof}
        $\underline{\Longrightarrow}$ Рассмотрим полилинейное отображение\\
        $L_T(dx^{i_1}, \ \ldots, \ dx^{i_p}, \ \partial x_{j_1}, \ \ldots, \ \partial x_{j_q})=T^{i_1,\ldots,i_p}_{j_1,\ldots,j_q}$ и полилинейное отображение\\
        $L_T'(dx^{i_1'}, \ \ldots, \ dx^{i_p'}, \ \partial x_{j_1'}, \ \ldots, \ \partial x_{j_q'})=T^{i_1',\ldots,i_p'}_{j_1',\ldots,j_q'}$.
        $$L_T'(dx^{i_1}, \ \ldots, \ dx^{i_p}, \ \partial x_{j_1}, \ \ldots, \ \partial x_{j_q})=L_T'\left(\frac{\partial x^{i_1}}{\partial x^{i_1'}}dx^{i_1'}, \ \ldots, \ \frac{\partial x^{i_p}}{\partial x^{i_p'}}dx^{i_p'}, \ \frac{\partial x^{j_1'}}{\partial x^{j_1}}\partial x_{j_1'}, \ \ldots, \ \frac{\partial x^{j_q'}}{\partial x^{j_q}}\partial x_{j_q'}\right)=$$
        $$=\frac{\partial x^{I_p}}{\partial x^{I_p'}}\cdot\frac{\partial x^{J_q'}}{\partial x^{J_q}}\cdot L_T'(dx^{i_1'}, \ \ldots, \ dx^{i_p'}, \ \partial x_{j_1'}, \ \ldots, \ \partial x_{j_q'})=\frac{\partial x^{I_p}}{\partial x^{I_p'}}\cdot\frac{\partial x^{J_q'}}{\partial x^{J_q}}\cdot T_{J_q'}^{I_p'}\overset{(1)}{=}T_{J_q}^{I_p}=$$
        $$=L_T(dx^{i_1}, \ \ldots, \ dx^{i_p}, \ \partial x_{j_1}, \ \ldots, \ \partial x_{j_q})$$
        Таким образом $L_T=L_T'$, то есть отображение задано однозначно.\\
        (1) По первому определению тензора
        $\underline{\Longleftarrow}$ Пусть $L:\left(T_PM\right)^p\times\left(T_P^*M\right)^q$ --- полилинейное отображение.\\
        Рассмотрим $(x^{1}, \ \ldots, \ x^{n})\longrightarrow T^{i_1,\ldots,i_p}_{j_1,\ldots,j_q}=L(dx^{i_1}, \ \ldots, \ dx^{i_p}, \partial x^{j_1}, \ \ldots, \ \partial x^{j_q})$,\\$(x^{1'}, \ \ldots, \ x^{n'})\longrightarrow T^{i_1',\ldots,i_p'}_{j_1',\ldots,j_q'}=L(dx^{i_1'}, \ \ldots, \ dx^{i_p'}, \partial x^{j_1'}, \ \ldots, \ \partial x^{j_q'})$
        $$T^{I_p}_{J_q}=L_T(dx^{i_1}, \ \ldots, \ dx^{i_p}, \ \partial x_{j_1}, \ \ldots, \ \partial x_{j_q})=$$
        $$=L_T\left(\frac{\partial x^{i_1}}{\partial x^{i_1'}}dx^{i_1'}, \ \ldots, \ \frac{\partial x^{i_p}}{\partial x^{i_p'}}dx^{i_p'}, \ \frac{\partial x^{j_1'}}{\partial x^{j_1}}\partial x_{j_1'}, \ \ldots, \ \frac{\partial x^{j_q'}}{\partial x^{j_q}}\partial x_{j_q'}\right)=$$
        $$=\frac{\partial x^{I_p}}{\partial x^{I_p'}}\cdot\frac{\partial x^{J_q'}}{\partial x^{J_q}}\cdot L_T(dx^{i_1'}, \ \ldots, \ dx^{i_p'}, \ \partial x_{j_1'}, \ \ldots, \ \partial x_{j_q'})=\frac{\partial x^{I_p}}{\partial x^{I_p'}}\cdot\frac{\partial x^{J_q'}}{\partial x^{J_q}}\cdot T_{J_q'}^{I_p'}$$
        Это равносильно $T_{J_q'}^{I_p'}= T_{J_q}^{I_p}\cdot\frac{\partial x^{I_p'}}{\partial x^{I_p}}\cdot\frac{\partial x^{J_q}}{\partial x^{J_q'}}$.
    \end{proof}
    \begin{lemma}\tab
        \begin{enumerate}
            \item dim$T_q^p=n^{p+q}$
            \item $T=T^{I_p}_{J_q}\cdot D^{J_q}_{I_p}$ и $D^{J_q}_{I_p}$ линейно независимы, где\\ $D^{J_q}_{I_p}(dx^{a_1}, \ \ldots, \ dx^{a_p}, \ \partial x^{b_1}, \ \ldots, \ x^{b_q})=\begin{cases}
                1, \ \textup{только на одном наборе}\\
                0, \ \textup{на остальных}
            \end{cases}$
        \end{enumerate}
    \end{lemma}
    \begin{proof}
        $T^{I_p}_{J_q}\cdot D^{J_q}_{I_p}(dx^{a_1}, \ \ldots, \ dx^{a_p}, \ \partial x^{b_1}, \ \ldots, \ x^{b_q})=T^{a_1,\ldots a_p}_{b_1,\ldots,b_q}$, так как все слагаемые, кроме одного, будут равны нулю. А так как слагаемых $n^{p+q}$, то dim$T^p_q=n^{p+q}$.
    \end{proof}
\end{document}